% Nous -- technical report (LaTeX rendering of docs/REPORT.md; keep in sync in content).
% Self-contained: compiles with pdflatex (no bibtex pass; uses thebibliography).
% If you have no local TeX toolchain, paste this file at https://overleaf.com to compile free.
% Figures are in ./figures/ (PNG, rendered from docs/diagrams/*.mmd).
\documentclass[11pt]{article}

\usepackage[utf8]{inputenc}
\usepackage[T1]{fontenc}
\usepackage{lmodern}
\usepackage[margin=1in]{geometry}
\usepackage[strings]{underscore} % print _ literally in text (e.g. test_nous_store.py)
\usepackage{amssymb}             % \geq \leq \in
\usepackage{graphicx}
\usepackage{booktabs}
\usepackage{float}
\usepackage{caption}
\usepackage{enumitem}
\usepackage{xcolor}
\usepackage[hidelinks]{hyperref}

\setlength{\parskip}{0.5em}
\setlength{\parindent}{0pt}

\newcommand{\code}[1]{\texttt{#1}}
% figure helper: caps width and height with keepaspectratio so any aspect fits a page
\newcommand{\nousfig}[3]{%
  \begin{figure}[htbp]\centering
    \includegraphics[width=\linewidth,height=0.78\textheight,keepaspectratio]{figures/#1}%
    \caption{#2}\label{#3}%
  \end{figure}}

\title{\textbf{Nous: Adversarial Testing of a Dependency-Free\\
Memory-and-Activation Substrate for Cognitive Architectures}\\[4pt]
\large An Engineering and Methodology Case Study}
\author{Cagan Caliskan}
\date{2026}

\begin{document}
\maketitle

\begin{abstract}
\noindent
This report describes \textbf{Nous}, a dependency-free Python prototype of a memory-and-activation
substrate for a cognitive architecture, and (more centrally) the \textbf{software-engineering
methodology} used to build it. Nous comprises four layers: a graph \emph{substrate store} of concept
nodes and weighted ``roads''; a deterministic \emph{spreading-activation} engine; a multi-source
\emph{activation field}; and a thin end-to-end \emph{perceive $\rightarrow$ activate $\rightarrow$ form
coalitions $\rightarrow$ select a thought} slice. The lower two layers are verified, by an adversarial
property-style test suite, for three concrete properties: \textbf{determinism} (identical output across
process runs, confirmed in fresh interpreters under randomized hash seeds), \textbf{integrity} (graph
indices stay mutually consistent across every mutation), and \textbf{purity} (reads never mutate stored
state). The upper two layers are \textbf{exploratory}: they run and are deterministic, but their
cognitive \emph{meaningfulness is explicitly not claimed}, and on clustered real input the end-to-end
slice produces an \emph{incoherent} ``thought,'' recorded as an asserted \textbf{negative result /
known limitation}, not a contribution. The central finding is a recurring failure mode: a component
that \emph{passes on a friendly, hand-chosen input and fails on an adversarial or system-generated
one}, and the two practices that caught it: \textbf{adversarial/property-based testing} (including
cross-process determinism and read-purity checks) and \textbf{loop-closure} (driving each stage with
system-generated rather than hand-placed inputs). Together they exposed \textbf{six distinct structural
bugs} that ordinary happy-path tests pass over. The artifact's verified totals: \textbf{130 assertions
across 56 test functions, all passing, in five suites.}
\end{abstract}

\section{Introduction}

Most write-ups of ambitious systems describe what the system is meant to become. This one separates
that aspiration from what was actually built and proven, and reports the engineering discipline in
between. The subject is a cognitive-architecture prototype; the contribution is the \textbf{testing
methodology} applied to it.

\textbf{Motivation:} Cognitive-architecture components (graph memories, spreading activation,
clustering into ``coalitions'') are easy to implement in a way that \emph{looks} correct on a
hand-built example and is quietly wrong in structure: it double-counts, it depends on iteration order,
it bakes the answer into the input. These are not typos; they are bugs in the \emph{shape} of the
algorithm, and they survive happy-path tests. The question this project answers in the affirmative is
whether a disciplined, adversarial, property-oriented test process can surface such bugs early and
force the honest scoping of what has and has not been achieved.

\textbf{Aspiration vs.\ built reality:} Nous began as a large design vision. Those original vision
documents are retained verbatim in the repository under \code{docs/design/}, but they are
\textbf{quarantined as aspiration}: a disclaimer there states plainly that their grand language
(``AGI,'' ``a real brain,'' ``mind,'' ``felt experience'') describes intent, not the built code, and
that of the thirteen documents only the two layer-plan files correspond to anything implemented. This
report concerns \textbf{only the built artifact} and is bound by the scope below.

\textbf{Scope and non-claims (binding):}
\begin{itemize}[nosep]
  \item Nous is an \textbf{engineering prototype} and a \textbf{testing-methodology case study}.
  \item It is \textbf{not} a mind, consciousness, AGI, or a validated cognitive model, and makes
        \textbf{no claim of machine cognition}.
  \item The substrate store and spreading-activation engine are \textbf{proven} deterministic,
        integrity-preserving, and pure. These are claims about \emph{mechanism}.
  \item The field, coalition, and ``thought'' layers are \textbf{exploratory}. Their cognitive
        meaningfulness is \textbf{not established}; the incoherent-thought behavior on clustered real
        input is an \textbf{honest negative result}, not a feature.
  \item Named neuroscience and cognitive theories are \textbf{conceptual inspiration} only. Nothing
        here \emph{implements} or \emph{reproduces} them, and no biological fidelity is claimed.
\end{itemize}

\textbf{Contributions (methodology first):}
\begin{enumerate}[nosep]
  \item A \textbf{software-engineering experience report} showing that adversarial/property-based
        testing (determinism verified across process hash seeds, and read-purity verified by hammering
        reads) and \textbf{loop-closure} caught a sequence of \textbf{six structural bugs} in a
        memory/activation substrate.
  \item A \textbf{verified substrate}: a graph store and a spreading-activation engine that are
        deterministic, integrity-preserving, and pure, with the verification shipped as runnable tests.
  \item An \textbf{honestly-scoped exploratory cognitive layer}, including an \emph{asserted}
        known-limitation test that encodes a negative result as a standing, checkable claim.
  \item A \textbf{reusable discipline}: transferable lessons (\S\ref{sec:lessons}) about how ``passes on
        a friendly input'' recurs and how to design tests and inputs that defeat it.
\end{enumerate}

\section{Background and related work}

The following bodies of work \textbf{inspired} the design's vocabulary and structure. Nous does not
implement, validate, or reproduce any of them, and claims no biological fidelity; each is named to
position the prototype, not to assert a result.

\textbf{Complementary Learning Systems (CLS):} The fast/slow split between a rapidly-learning episodic
system and a slowly-consolidating semantic one~\cite{mcclelland1995}, updated for intelligent-agent
design in~\cite{kumaran2016}, \emph{inspires} Nous's separation of \textbf{activation} (what lights up
now) from \textbf{consolidation} (what is later reinforced). In the prototype this appears only as a
strict rule: spreading activation is a pure read and never strengthens what it touches, not as a
consolidation system.

\textbf{Synaptic Tagging and Capture (STC):} The hypothesis that a weak event sets a local ``tag'' that
can capture plasticity-related products from a nearby strong event~\cite{frey1997, redondo2011}
\emph{inspires} the framing of reinforcement as a \emph{gated, budgeted} event, which is why Nous keeps
activation read-only; it does not model tagging dynamics.

\textbf{Global Workspace Theory (GWT):} The idea that specialized processes form coalitions that
compete, and a winner is broadcast globally~\cite{baars1988, dehaene2011}, is the source of Nous's
\textbf{coalition} vocabulary. Only coalition \emph{formation} exists in the code; the
competition/broadcast step is an explicit, unbuilt placeholder.

\textbf{LIDA cognitive cycle:} LIDA's systems-level \emph{perceive $\rightarrow$ attend $\rightarrow$
broadcast} cycle~\cite{franklin2014} loosely shapes the staging of the end-to-end slice, which realizes
only the early stages.

\textbf{Cognitive architectures:} ACT-R's integrated theory of cognition~\cite{anderson2004} and
SPAUN's large-scale functional brain model~\cite{eliasmith2012} are referenced for \textbf{positioning}:
Nous is far smaller and makes no behavioral or biological claims; it is a substrate-and-method study,
not an architecture in their sense.

\textbf{Spreading-activation theory:} Nous's activation engine is a deterministic, engineered variant of
spreading activation in the sense of Collins \& Loftus~\cite{collins1975}; the convergence-sums property
is from that lineage, while the hop-layered single-expansion is an engineering choice made for
determinism (\S\ref{sec:findings}).

\textbf{Arousal and plasticity:} The Yerkes--Dodson inverted-U relating arousal to
performance~\cite{yerkes1908} appears in the design notes as motivation for emotion-modulated
plasticity; in the code this is only charge-modulated strengthening and is best read as design-level
aspiration.

\textbf{Semantic priming}~\cite{collins1975, meyer1971} is named here because it is the worked example
for future validation (\S\ref{sec:future}): an effect whose direction is known independently of Nous.

\section{System architecture (as built)}\label{sec:arch}

Nous is pure-Python standard library, organized bottom-up; each layer reads the one below through its
public surface only. Figure~\ref{fig:archstack} makes the honest boundary visual; the figures that
follow detail the data model, the read/write paths, and the end-to-end slice.

\nousfig{architecture_stack.png}{\textbf{The four-layer stack and the honest boundary:} Layers 1--2
(substrate store, spreading activation) are PROVEN (determinism, integrity, purity); layers 3--4
(multi-source field, end-to-end slice) are EXPLORATORY. It establishes \emph{where verification stops};
it does \emph{not} claim the exploratory layers are cognitively meaningful.}{fig:archstack}

\subsection{Substrate store (\code{nous_store.py}, 548 lines)}
The only component that mutates the graph; it holds no cognitive logic.
\begin{itemize}[nosep]
  \item \textbf{Nodes and roads:} A \code{Node} is an always-neutral unit of information with a stable
        identity and a content string. A \code{Road} is a \emph{directed} edge $\text{source}
        \rightarrow \text{target}$ carrying a scalar \code{strength} (a new road starts at $0.05$) and
        a \emph{multi-valued} emotional life \code{charges} (named feelings $\rightarrow$ intensity in
        $[-1,1]$), from which an aggregate \code{valence} is derived. There is \textbf{one road per
        $(\text{source}, \text{target})$ pair}; the reverse direction is a distinct road.
  \item \textbf{Stable identity vs.\ ordering:} Every node and road has a stable \code{uuid4} identity
        and is never addressed by position; a stale handle resolves to \code{NotFound}, never to a wrong
        object. Formation order, when needed, is a \textbf{separate integer \code{creation_seq}}, kept
        distinct from identity precisely so ordering never has to fall back to a uuid
        (\S\ref{sec:findD}).
  \item \textbf{Lazy, pure-read decay:} Roads decay with elapsed time, but decay is \textbf{computed on
        access as a read-only view}: a read returns a \emph{copy} carrying decayed values and never
        writes. The strength time-constant scales steeply with strength
        ($\tau = 45 \cdot \exp(\text{strength}\cdot 4.2)$), so a ``highway'' is far more stable than a
        ``dirt track.'' This purity is load-bearing (finding~(a)).
  \item \textbf{Sticky closure:} A strongly-negative crossing (or one-shot trauma) \textbf{closes} a
        road; closure is re-evaluated \emph{only on an active crossing}, never by a read or by the
        passage of time, and only an active positive crossing reopens it.
  \item \textbf{Pin guard and integrity:} \code{remove_node} refuses to delete a \emph{pinned} node;
        adjacency, dedup, and designation consistency are asserted incrementally on every mutation,
        with a full-graph scan available offline.
  \item \textbf{Atomicity:} Every mutation is a synchronous method with no \code{await} inside, making
        each mutation atomic by construction under single-threaded asyncio.
\end{itemize}

\nousfig{data_model.png}{\textbf{Substrate data model:} The \code{Node} and \code{Road} fields;
identity (\code{uuid4}) is deliberately separate from ordering (\code{creation_seq}). Describes the
data structures only.}{fig:datamodel}

\nousfig{substrate_read_write.png}{\textbf{Read vs.\ write paths (the purity invariant):} Reads compute
decay into a \emph{copy} and never mutate; only writes persist decay; closure changes only on an active
crossing. Visualizes finding (a); it proves a mechanism property, not anything about cognition.}{fig:readwrite}

\subsection{Spreading activation (\code{nous_activation.py}, 281 lines)}
The first operation that reads \emph{across} the graph: radiate activation from a seed along directed
out-roads and report what lit up. \textbf{Pure read:}
\begin{itemize}[nosep]
  \item \textbf{Convergence sums:} Activation arriving at a node from several paths is summed.
  \item \textbf{Hop-layered single-expansion:} The wave expands in hop layers; all hop-$k$
        contributions accumulate before any hop-$k$ node expands, each node expands \textbf{exactly
        once}, and a per-layer \textbf{frozen snapshot} is used so same-hop (``intra-layer'') edges
        cannot leak iteration order into the result. Equal-length convergence is \textbf{exact}, and
        the algorithm is a deterministic function of store state. This replaced an earlier
        priority-frontier design (findings~(b),~(c)).
  \item \textbf{Two attenuators:} Transmitted activation is
        $\text{outgoing}\cdot\text{road.strength}\cdot\text{decay\_per\_hop}$: \code{road.strength} is
        \textbf{edge conductance}; \code{decay_per_hop} is a \textbf{global distance penalty}. Kept
        separate deliberately.
  \item \textbf{Leaky closure:} A single swappable policy governs closed roads: below a breakthrough
        threshold nothing crosses; above it, a small attenuated fraction leaks through and is logged.
        The constants are \textbf{uncalibrated guesses on probation}; the tests pin only the qualitative
        regime.
  \item \textbf{Flow recording:} The engine records, per edge, the activation actually transmitted
        (\code{flow}), distinct from static topology. Downstream consumers read \code{flow}, not raw
        edges.
  \item \textbf{Locality:} A hop budget and an activation floor keep the spread in the seed's
        neighborhood, which makes it affordable in pure Python.
\end{itemize}

\subsection{Multi-source activation field (\code{nous_field.py}, 161 lines)}
Runs the proven single-source spread once \textbf{per source} and assembles a field: each node's
activation from each source, the \textbf{summed total} across sources, and the \textbf{dominant source}
(the node's \emph{basin}). Dominant-source ties break on the stable \code{creation_seq}, so the field is
deterministic given store state. Pure read. It exists because a \emph{single} seed yields one activation
peak, so any segmentation of it keys off distance from the seed rather than community structure
(finding~(e)); multiple sources yield multiple genuine peaks.

\subsection{Coalitions and the end-to-end slice (\code{nous_coalitions.py} 324; \code{nous_perceive.py} 162)}
\begin{itemize}[nosep]
  \item \textbf{Coalition formation} groups co-activated nodes via a \emph{flow-cut} segmentation of a
        single spread or a \emph{group-by-basin} path for the multi-source field. Every
        output-affecting order is an explicit total order over intrinsic keys (activation, then
        \code{creation_seq}), never set-iteration order.
  \item \textbf{Perception} turns text into sources: tokens (length $\geq 3$) are matched by substring
        against node content, and matched nodes become activation sources with energy scaled by match
        count. Minimal: no embeddings, but it makes sources \textbf{input-derived rather than
        hand-placed}.
  \item \textbf{The slice} (\code{think}) runs perceive $\rightarrow$ build field $\rightarrow$ form
        coalitions by basin $\rightarrow$ select the strongest coalition as ``the thought.'' The final
        step is an explicit \textbf{placeholder} for an unbuilt workspace-competition/broadcast phase.
        The slice is pure and deterministic; what it produces is \textbf{not} claimed to be meaningful
        (finding~(f), \S\ref{sec:negresult}).
\end{itemize}

\nousfig{pipeline.png}{\textbf{End-to-end ``thought'' slice:} text $\rightarrow$ perceive $\rightarrow$
field $\rightarrow$ coalitions $\rightarrow$ thought. The final ``strongest coalition = thought'' step is
a placeholder for an unbuilt broadcast phase; the figure shows \emph{data flow}, not that the resulting
``thought'' is meaningful.}{fig:pipeline}

\section{Methodology}\label{sec:method}

The central claim of this report is that the test process (not the prototype) is the contribution. Two
properties were treated as first-class and verified directly.

\textbf{Property-based / adversarial testing:} Each component ships with tests written to \emph{break}
its central claim, not to confirm it on an easy case: lopsided edge strengths, intra-layer edges,
equal-activation rings, and a two-dense-communities-joined-by-a-thin-bridge graph built specifically to
defeat distance-based segmentation.

\textbf{Determinism across processes, not within one:} The natural implementations leak \code{set}-
iteration order over \code{uuid4} keys into their results. Within a single process the hash seed is
fixed, so this cannot be seen; the determinism tests therefore spawn \textbf{fresh interpreters} with
different \code{PYTHONHASHSEED} values and compare output byte-for-byte. The activation suite uses seeds
$\{0,1,3,7,42\}$; the coalition and field suites use $\{0,1,2,3,7,42,123,999\}$ (eight seeds); the slice
uses $\{0,1,2,7,42,999\}$.

\textbf{Read-purity as an asserted invariant:} Read-only operations are run in tight loops and the full
canonical (un-decayed) store state is asserted byte-equal before and after, in every layer.

\textbf{The harness, and a deliberate non-use of pytest:} The suites use a tiny built-in harness: each
\code{check(...)}\,/\,\code{expect_raises(...)} increments a counter, and the suite prints
\code{RESULT: N passed, M failed} and exits non-zero on any failure. They are intentionally \textbf{not}
pytest tests: pytest would collect the \code{test_*} functions but would not register a logical
\code{check(False, ...)} as a failed assertion, so it could report green on a real failure. The honest
runner is \code{run_tests.py}, which sets an absolute \code{PYTHONPATH} (required so the determinism
subprocesses resolve their imports regardless of working directory) and propagates non-zero exit codes.

\section{Findings overview}

Read findings (b)--(f) together and one shape repeats: \textbf{a component passes on a friendly,
hand-chosen input and fails on an adversarial or system-generated one.} Each fix had two parts: write
the test that \emph{tries to break} the claim, and make the inputs \emph{system-generated} rather than
hand-chosen so a stage can no longer be validated against a fiction of what feeds it. The six findings
below are each presented as: \emph{claim under test $\rightarrow$ bug $\rightarrow$ root cause
$\rightarrow$ fix $\rightarrow$ the specific test that now locks it.} Every claim is traceable to the
code and the named test; reported magnitudes that the shipped artifact does not let one reproduce are
omitted rather than asserted.

\nousfig{friendly_input_loop.png}{\textbf{The recurring ``friendly-input'' failure mode, and the two
remedies.} Conceptual: the loop shared by findings (b)--(f), broken by adversarial testing and by closing
the loop. It illustrates the \emph{method}; it asserts no measurement.}{fig:friendly}

\section{The six structural findings}\label{sec:findings}

\nousfig{six_findings.png}{\textbf{The six structural findings at a glance:} Each row: claim
$\rightarrow$ bug $\rightarrow$ root cause $\rightarrow$ fix $\rightarrow$ locking test. Every cell traces
to a module docstring and a named test; no error magnitudes are asserted.}{fig:sixfindings}

\subsection{(a) Read-as-write: a read that could silently re-open a closed wound}
\textbf{Claim under test:} reads are pure; stored state changes only on writes.
\textbf{Bug:} the natural implementation applies time-decay on access and writes it back, so every read
mutates the store; worse, if decay also re-evaluated closure, repeated reads over time could let a
\emph{closed} (traumatic) edge drift back open with no crossing ever reopening it.
\textbf{Root cause:} conflating ``compute the current value'' with ``persist the current value,'' and
letting a passive read touch state (\code{closed}) that should change only on an active event.
\textbf{Fix:} split decay into a pure view (\code{_compute_decay} returns decayed values without
mutating; \code{get_road} returns a \emph{copy}) and a persisting write (\code{_persist_decay}, called
only by writes/offline). Decay never touches \code{closed}; closure is sticky.
\textbf{Locked by:} \code{test_reads_are_pure}; \code{test_lazy_decay}; \code{test_closure_is_sticky}
(time + repeated reads do not reopen a closed road; only an active positive return does).

\nousfig{sticky_closure_state.png}{\textbf{A road's closure lifecycle (finding (a) invariant):}
Open/Closed transitions happen \emph{only} on active crossings; time passing and reads never reopen a
closed road. A mechanism property, asserted by \code{test_closure_is_sticky}.}{fig:closure}

\subsection{(b) Convergence double-counting: ``first-pop-final'' where activation sums}
\textbf{Claim under test:} a node downstream of a convergence point is not over-counted.
\textbf{Bug:} when activation converges from several paths it must sum, but re-expanding the convergence
node as later paths raise its value counts everything downstream multiple times.
\textbf{Root cause:} importing Dijkstra's ``first-pop-final'' intuition into an algorithm where
activation \emph{accumulates} rather than being a shortest-path minimum.
\textbf{Fix:} bounded single-expansion, hop-layered: each node expands exactly once carrying the full
sum of paths of length $\leq$ its settle hop; equal-length convergence is exact.
\textbf{Locked by:} \code{test_no_double_count_downstream_of_convergence}, which uses \emph{deliberately
lopsided} edge strengths to prove the result does not depend on them and includes an explicit guard that
the result is \emph{strictly below} what additive re-expansion would produce;
\code{test_convergent_activation_sums}.

\nousfig{convergence_before_after.png}{\textbf{Convergence double-counting, before/after (finding (b)):}
Structural schematic of the diamond-with-tail: naive re-expansion over-counts everything downstream of
the convergence node; single-expansion counts it once. No magnitudes are shown (the report omits them as
not reproducible).}{fig:convergence}

\subsection{(c) Order-dependence, then nondeterminism: a priority frontier over hashed sets}
\textbf{Claim under test:} post-convergence activation depends only on the graph, never on expansion
order or process hash seed.
\textbf{Bug:} an activation-priority frontier made \emph{which} late ripple was dropped depend on
incidental edge-strength ordering; and with same-hop edges, which sibling expanded first was decided by
iterating a \code{set} of uuids: order randomized per process by \code{PYTHONHASHSEED}, so results
swung across runs.
\textbf{Root cause:} letting expansion order be load-bearing for the result, over a container whose
iteration order is intentionally randomized.
\textbf{Fix:} hop-layered expansion with \textbf{per-layer frozen values}: a layer's expansion snapshot
is fixed before any node in it expands, so an intra-layer edge's contribution propagates only in the
next layer.
\textbf{Locked by:} \code{test_expansion_is_deterministic_under_edge_strength}, and the decisive
\code{test_determinism_with_same_hop_edge_adversarial}, which builds an intra-layer edge and re-runs the
spread in fresh interpreters under $\text{PYTHONHASHSEED} \in \{0,1,3,7,42\}$, asserting byte-identical
results.

\subsection{(d) Identity/order leak: tie-breaking on a uuid regenerated every run}\label{sec:findD}
\textbf{Claim under test:} coalition output (partition, ordering, representatives) is identical across
processes.
\textbf{Bug:} coalition ordering and representative selection tie-broke on the node id, but ids are
\code{uuid4}, regenerated every run, so any tiebreak falling back to the id was nondeterministic across
processes even when the structure was identical.
\textbf{Root cause:} using identity (intentionally non-stable across runs) to carry ordering.
\textbf{Fix:} tie-break on the stable \code{creation_seq} integer; \code{representative()} orders by
$(-\text{activation}, \text{creation\_seq})$ and the coalition sort key is
$(-\text{strength}, -\text{size}, \min \text{creation\_seq})$.
\textbf{Locked by:} \code{test_determinism_adversarial_across_hash_seeds}, which builds a ring of
equal-activation nodes plus tie-strength singletons, serializes by stable keys (content +
\code{creation_seq}, never uuid), and asserts identity across 8 hash seeds;
\code{test_representative_is_most_active_deterministically}.

\subsection{(e) Binding degeneracy: components on a graph connected by construction}
\textbf{Claim under test:} coalition binding subdivides a co-activated region into the right groups.
\textbf{Bug, two stages:} (1) Connectivity binding is structurally degenerate: a seed-rooted spread's
co-activated subgraph is \textbf{connected by construction}, so connected-components always yields
\emph{one} coalition. (2) Flow-magnitude cutting does subdivide, but for a single seed the only strong
signal is \textbf{distance from the seed}, and flow magnitude is even \emph{anti-correlated} with
community membership near the seed.
\textbf{Root cause:} a single seed produces one activation peak and monotonically-decaying flow; any
segmentation of it keys off distance, not community.
\textbf{Fix:} keep flow-cut as the single-seed segmenter (it does yield more than one coalition) but
recognize the real problem is \emph{upstream} and reframe to a \textbf{multi-source field}: multiple
sources $\rightarrow$ multiple genuine peaks $\rightarrow$ basins around peaks.
\textbf{Locked by:} \code{test_single_seed_spread_yields_multiple_coalitions};
\code{test_acceptance_multisource_recovers_communities} (the field recovers exactly the two communities
on the adversarial bridge graph that single-seed segmentation could not).

\nousfig{single_seed_vs_multisource.png}{\textbf{Single-seed degeneracy vs.\ multi-source basins
(finding (e)).} The same two-communities-with-a-thin-bridge graph: one seed yields a single
distance-keyed blob; two sources yield two basins with the bridge as the seam. The multi-source recovery
is the asserted acceptance test: in which the two sources are \emph{placed one per community} (see
finding (f) for why that matters).}{fig:singlemulti}

\subsection{(f) Circular validation: the answer hand-encoded in the input}\label{sec:findF}
\textbf{Claim under test:} the cognitive layer works on input the \emph{system} produced, not input
chosen to make it work.
\textbf{Bug:} every layer above the substrate had been validated against a fiction of what would feed
it. Sources were \textbf{hand-placed}, and this culminated in the field acceptance test recovering two
communities \emph{only because one source was hand-placed at the center of each}: the answer was
encoded in the input. (This is visible in the artifact:
\code{test_acceptance_multisource_recovers_communities} does place one source per community by hand. It
is honest as a \emph{mechanism} test and would be dishonest if read as evidence that the system finds
communities on its own.)
\textbf{Root cause:} validating a stage against inputs invented to produce the desired output.
\textbf{Fix:} \textbf{close the loop} with the thinnest real slice (\code{nous_perceive.py}), where which
nodes become sources is decided by the input matched against the substrate, not by the test. On real
clustered input the field then produces an \textbf{incoherent} ``thought'': basins collapse, one source
over-swallows across domains, and same-domain sibling sources strand as singletons.
\textbf{Locked by: as a \emph{limitation}, not a result:}
\code{test_KNOWN_LIMITATION_clustered_sources_give_poor_basins} \emph{asserts the failure on purpose} and
is written to flip to green only when source-grouping is fixed.

\subsection{What the six have in common}
(b)--(f) are all instances of ``passes on a friendly input.'' (a) is the dual lesson for
\emph{invariants}: a property (purity) that an easy test trivially satisfies can be violated by an
interaction (decay $\times$ closure $\times$ time) that only a purpose-built test exercises. In every
case the bug was structural and the catching test was adversarial or loop-closing.

\nousfig{cognitive_evolution.png}{\textbf{The cognitive-layer arc and its dead-ends:} connectivity
binding $\rightarrow$ flow-cut $\rightarrow$ multi-source field $\rightarrow$ closed-loop, and what each
step revealed (findings (e), (f)), ending at the known limitation. It documents the \emph{development
path}, not a working capability.}{fig:evolution}

\subsection{Transferable lessons}\label{sec:lessons}
\begin{enumerate}[nosep]
  \item \textbf{Test the property, not the example:} assert purity, determinism, and ``no
        double-count'' directly and across processes, not by spot-checking outputs.
  \item \textbf{Randomize what you assume is fixed:} running under multiple \code{PYTHONHASHSEED}
        values is what makes uuid-set nondeterminism observable.
  \item \textbf{Close the loop:} let a stage's inputs be produced by the stage below it; hand-chosen
        inputs hide both degeneracy and circularity.
  \item \textbf{Encode known limitations as asserted tests:} a negative result written as a
        failing-on-purpose assertion is a standing, checkable claim, not a comment that rots.
\end{enumerate}

\section{Results}

\subsection{Verified properties (the positive results)}
For the substrate store and the spreading-activation engine, the test suite establishes, as runnable
checks: \textbf{determinism} (identical output across process runs, confirmed in fresh interpreters under
the hash seeds in \S\ref{sec:method}); \textbf{integrity} (graph indices stay mutually consistent across
every mutation; a full scan finds no drift); and \textbf{purity} (reads never mutate stored state). These
are claims about \emph{mechanism}: the code does exactly and only what it says, reproducibly, not that
the mechanism reproduces any cognitive phenomenon.

\nousfig{determinism_seeds.png}{\textbf{Adversarial determinism across process hash seeds
(before/after).} \emph{AFTER} is the asserted, byte-identical result of the four determinism tests (each
across its own seed set); \emph{BEFORE} is the documented prior bug (findings (c), (d)), shown
qualitatively and \emph{not} re-measured. Reproducibility of mechanism only, not cognition.}{fig:determinism}

\subsection{Test suite}
All figures below are taken verbatim from \code{docs/TEST_OUTPUT.txt} (CPython 3.13.2, Darwin 24.5.0
arm64) and re-confirmed by \code{python3 run_tests.py} for this report. Every suite passes with zero
failures.

\begin{table}[H]
\centering
\begin{tabular}{llr}
\toprule
\textbf{Suite} & \textbf{Layer} & \textbf{Checks passed} \\
\midrule
\code{test_nous_store.py}      & substrate store               & 45 \\
\code{test_nous_activation.py} & spreading activation          & 35 \\
\code{test_nous_coalitions.py} & coalition formation           & 19 \\
\code{test_nous_field.py}      & multi-source field            & 13 \\
\code{test_nous_perceive.py}   & perception + end-to-end slice & 18 \\
\midrule
\textbf{Total} &               & \textbf{130 passed, 0 failed} \\
\bottomrule
\end{tabular}
\caption{Per-suite checks (verbatim from \code{docs/TEST_OUTPUT.txt}).}
\end{table}

The total is \textbf{130 assertions executed across 56 test functions}, including the cross-process
determinism checks and one explicitly-asserted known-limitation test. The verbatim run is reproduced in
Appendix~A.

\begin{table}[H]
\centering
\begin{tabular}{lrrr}
\toprule
 & \textbf{Files} & \textbf{Lines (\code{wc -l})} & \textbf{Non-blank, non-comment} \\
\midrule
Source modules (\code{src/}) & 5  & 1{,}476 & 1{,}058 \\
Test suites (\code{tests/})  & 5  & 1{,}555 & 1{,}223 \\
\midrule
\textbf{Python total}        & \textbf{10} & \textbf{3{,}031} & \textbf{2{,}281} \\
\bottomrule
\end{tabular}
\caption{Measured code size. Per module (\code{wc -l}): \code{nous_store.py} 548,
\code{nous_activation.py} 281, \code{nous_coalitions.py} 324, \code{nous_field.py} 161,
\code{nous_perceive.py} 162.}
\end{table}

\nousfig{test_suite_bars.png}{\textbf{Passing assertions per suite (130 total, 0 failed):} Counts taken
verbatim from \code{docs/TEST_OUTPUT.txt}. Proves \emph{mechanism} (and the \emph{mechanics} of the
exploratory layers); it does \emph{not} measure cognitive quality. Colouring matches
Figure~\ref{fig:archstack}.}{fig:bars}

\subsection{The incoherent-thought known limitation (an asserted negative result)}\label{sec:negresult}
On clustered real input the end-to-end slice does \textbf{not} produce coherent coalitions. When
perception lights up several nodes in the same domain, basins collapse: one source over-swallows across
domains while same-domain siblings strand as singletons. This is asserted on purpose by
\code{test_KNOWN_LIMITATION_clustered_sources_give_poor_basins}, which passes today \emph{because the
failure holds}; it is written to flip only when the underlying source-grouping problem is solved. It is
reported here as a \textbf{negative result}, not a contribution.

\nousfig{closed_loop_known_limitation.png}{\textbf{The closed-loop slice and its honest failure
(finding (f); asserted NEGATIVE result).} On the real input ``coffee in the kitchen this morning,'' the
``thought'' over-swallows across domains (pulls in \emph{work}/\emph{code}) and same-domain perceived
sources strand as singleton basins, asserted by
\code{test_KNOWN_LIMITATION_clustered_sources_give_poor_basins}. A documented \textbf{known limitation},
not a capability.}{fig:closedloop}

\subsection{What these results do and do not establish}
They establish that the substrate and activation engine are correct, deterministic,
integrity-preserving, and pure, and that the cognitive layers run deterministically and purely. They do
\textbf{not} establish that any ``thought'' the cognitive layers produce is cognitively meaningful; on
the contrary, the one end-to-end behavioral probe on real input is a documented failure. ``The mechanism
is correct and deterministic'' and ``the mechanism reproduces a cognitive phenomenon'' are different
claims; this work earns the first and explicitly does not claim the second.

\section{Limitations and threats to validity}
\begin{itemize}[nosep]
  \item \textbf{Cognitive meaningfulness is unestablished:} The field/coalition/thought layers are
        exploratory; the only end-to-end behavioral result on real input is the negative one above.
  \item \textbf{Uncalibrated constants on probation:} The closed-road breakthrough threshold/leak
        fraction and the coalition flow-cut fraction are first-pass guesses; tests pin only the
        qualitative regime, and nothing downstream may depend on the magnitudes.
  \item \textbf{Perception is exact/substring matching only:} No embeddings or fuzzy semantics.
  \item \textbf{Out-roads-only spread (v1):} Backward spread along in-roads is a deliberate later
        addition.
  \item \textbf{Determinism is ``given store state,'' not over the abstract graph:} The proven property
        is that output is a deterministic function of the \emph{store's state}, where state includes the
        integer \code{creation_seq} assigned in build order. Two stores isomorphic as abstract graphs but
        built in a different order can differ in \code{creation_seq} and, where ties break on it, can
        produce different (but each individually deterministic and reproducible) tie-break orderings.
        Determinism here means reproducibility across process runs of the same construction, not
        invariance under graph isomorphism.
  \item \textbf{Scale is unverified:} All tests are small graphs; the prototype makes no scaling claims.
\end{itemize}

\section{Future work}\label{sec:future}
The natural next step is to stop validating the cognitive layer against hand-built inputs and instead
validate it against a \textbf{measurable cognitive effect whose direction is known independently of
Nous}. \textbf{Semantic priming} is the worked example: in human studies, recognizing a target word is
faster when preceded by a related prime than an unrelated one~\cite{meyer1971}, which the
spreading-activation account explains via pre-activation spreading to related concepts~\cite{collins1975}.
A priming-style protocol for Nous would pre-activate a related concept and test whether the target's
activation is measurably raised relative to an unrelated control: a directional prediction made
\emph{before} running the system, so the input is no longer chosen to make the answer come out right.
Two further items follow directly from the findings:
\begin{itemize}[nosep]
  \item \textbf{The source-grouping problem (from finding~(f)):} Same-domain perceived sources must be
        grouped or merged before basins are meaningful; solving it is the condition under which the
        known-limitation test would flip to green.
  \item \textbf{A real broadcast/competition phase:} The ``strongest coalition = thought'' placeholder
        should be replaced by an actual workspace competition, at which point the uncalibrated constants
        can begin to be calibrated against a running cycle rather than guessed.
\end{itemize}

\section{Conclusion}
What was earned is concrete and modest: a correct, deterministic, integrity-preserving, and pure
memory-and-activation substrate, built in dependency-free Python and shipped with the tests that prove
those properties across process hash seeds, and, above it, an honestly-scoped exploratory layer whose
one end-to-end behavioral probe is a documented negative result. The larger contribution is the
\textbf{method}: adversarial/property-based testing and loop-closure caught six structural bugs that
happy-path testing passes over, and the same discipline drew the boundary the project refuses to
cross: proving mechanism while explicitly declining to claim cognition. What remains open is exactly
what the method recommends next: evaluate the cognitive layer against an effect whose direction is known
independently, fix the source-grouping degeneracy, and build the competition phase that would make the
constants calibratable.

\begin{thebibliography}{99}
\bibitem{mcclelland1995} McClelland, J. L., McNaughton, B. L., \& O'Reilly, R. C. (1995). Why there are
complementary learning systems in the hippocampus and neocortex: Insights from the successes and failures
of connectionist models of learning and memory. \emph{Psychological Review, 102}(3), 419--457.
https://doi.org/10.1037/0033-295X.102.3.419
\bibitem{kumaran2016} Kumaran, D., Hassabis, D., \& McClelland, J. L. (2016). What learning systems do
intelligent agents need? Complementary Learning Systems theory updated. \emph{Trends in Cognitive
Sciences, 20}(7), 512--534. https://doi.org/10.1016/j.tics.2016.05.004
\bibitem{frey1997} Frey, U., \& Morris, R. G. M. (1997). Synaptic tagging and long-term potentiation.
\emph{Nature, 385}(6616), 533--536. https://doi.org/10.1038/385533a0
\bibitem{redondo2011} Redondo, R. L., \& Morris, R. G. M. (2011). Making memories last: The synaptic
tagging and capture hypothesis. \emph{Nature Reviews Neuroscience, 12}(1), 17--30.
https://doi.org/10.1038/nrn2963
\bibitem{baars1988} Baars, B. J. (1988). \emph{A cognitive theory of consciousness.} Cambridge University
Press. ISBN 0-521-30133-5.
\bibitem{dehaene2011} Dehaene, S., \& Changeux, J.-P. (2011). Experimental and theoretical approaches to
conscious processing. \emph{Neuron, 70}(2), 200--227. https://doi.org/10.1016/j.neuron.2011.03.018
\bibitem{franklin2014} Franklin, S., Madl, T., D'Mello, S., \& Snaider, J. (2014). LIDA: A systems-level
architecture for cognition, emotion, and learning. \emph{IEEE Transactions on Autonomous Mental
Development, 6}(1), 19--41. https://doi.org/10.1109/TAMD.2013.2277589
\bibitem{anderson2004} Anderson, J. R., Bothell, D., Byrne, M. D., Douglass, S., Lebiere, C., \& Qin, Y.
(2004). An integrated theory of the mind. \emph{Psychological Review, 111}(4), 1036--1060.
https://doi.org/10.1037/0033-295X.111.4.1036
\bibitem{eliasmith2012} Eliasmith, C., Stewart, T. C., Choo, X., Bekolay, T., DeWolf, T., Tang, Y., \&
Rasmussen, D. (2012). A large-scale model of the functioning brain. \emph{Science, 338}(6111),
1202--1205. https://doi.org/10.1126/science.1225266
\bibitem{collins1975} Collins, A. M., \& Loftus, E. F. (1975). A spreading-activation theory of semantic
processing. \emph{Psychological Review, 82}(6), 407--428. https://doi.org/10.1037/0033-295X.82.6.407
\bibitem{meyer1971} Meyer, D. E., \& Schvaneveldt, R. W. (1971). Facilitation in recognizing pairs of
words: Evidence of a dependence between retrieval operations. \emph{Journal of Experimental Psychology,
90}(2), 227--234. https://doi.org/10.1037/h0031564
\bibitem{yerkes1908} Yerkes, R. M., \& Dodson, J. D. (1908). The relation of strength of stimulus to
rapidity of habit-formation. \emph{Journal of Comparative Neurology and Psychology, 18}(5), 459--482.
https://doi.org/10.1002/cne.920180503
\end{thebibliography}

\appendix

\section{Verbatim test output}
The captured run from \code{docs/TEST_OUTPUT.txt} (per-suite result lines; complete capture, including
every check label and the \code{KNOWN LIMITATION} note, is that file in the repository).

\begin{verbatim}
NOUS SUBSTRATE STORE -- TEST SUITE (Layer 1)     -> RESULT: 45 passed, 0 failed  [exit 0]
NOUS SPREADING ACTIVATION -- TEST SUITE          -> RESULT: 35 passed, 0 failed  [exit 0]
NOUS COALITION FORMATION -- TEST SUITE           -> RESULT: 19 passed, 0 failed  [exit 0]
NOUS MULTI-SOURCE FIELD + FIELD COALITIONS       -> RESULT: 13 passed, 0 failed  [exit 0]
NOUS PERCEPTION + END-TO-END SLICE -- TEST SUITE  -> RESULT: 18 passed, 0 failed  [exit 0]
                                                    TOTAL: 130 passed, 0 failed
\end{verbatim}

\section{Reproduction}
\begin{verbatim}
# all five suites (handles the absolute PYTHONPATH for the determinism subprocesses):
python3 run_tests.py

# a single suite directly:
PYTHONPATH="$PWD/src" python3 tests/test_nous_store.py
\end{verbatim}
Pure standard library; no dependencies (developed and captured on CPython 3.13.2). The suites use a
built-in harness and must be run as scripts, not under pytest.

\section{Figure index and sources}
All figures are embedded inline in the sections above. Their editable sources and
presentation-resolution renders live in the repository: Mermaid diagrams as \code{docs/diagrams/*.mmd}
(rendered to \code{docs/diagrams/rendered/*.png}; the \code{.mmd} files also render crisply at
mermaid.live or via \code{mmdc}), and data charts as \code{docs/diagrams/figures/*.py} (matplotlib;
run them to regenerate the PNGs at \code{dpi=200}).

\begin{table}[H]
\centering\small
\begin{tabular}{rll l}
\toprule
\textbf{Fig} & \textbf{Section} & \textbf{Asset (in \code{docs/diagrams/})} & \textbf{Kind} \\
\midrule
1  & \S3   & \code{architecture_stack}            & Mermaid \\
2  & \S3.1 & \code{data_model}                    & Mermaid \\
3  & \S3.1 & \code{substrate_read_write}          & Mermaid \\
4  & \S3.4 & \code{pipeline}                      & Mermaid \\
5  & \S5   & \code{friendly_input_loop}           & Mermaid \\
6  & \S6   & \code{six_findings}                  & Mermaid \\
7  & \S6.1 & \code{sticky_closure_state}          & Mermaid \\
8  & \S6.2 & \code{convergence_before_after}      & Mermaid \\
9  & \S6.5 & \code{single_seed_vs_multisource}    & Mermaid \\
10 & \S6.7 & \code{cognitive_evolution}           & Mermaid \\
11 & \S7.1 & \code{figures/determinism_seeds.py}  & matplotlib \\
12 & \S7.2 & \code{figures/test_suite_bars.py}    & matplotlib \\
13 & \S7.3 & \code{closed_loop_known_limitation}  & Mermaid \\
14 & this appendix & \code{repo_map}              & Mermaid \\
\bottomrule
\end{tabular}
\end{table}

\nousfig{repo_map.png}{\textbf{Repository / asset map:} Orientation only; it shows the package layout,
not behavior.}{fig:repomap}

\end{document}
